\documentclass[12pt]{article}
\usepackage[utf8]{inputenc}
\usepackage[T1]{fontenc}
\usepackage[margin=1in]{geometry}
\usepackage{setspace}
\usepackage{xcolor}
\usepackage[colorlinks=true, urlcolor=blue!60!black, linkcolor=blue!60!black]{hyperref}
\setlength{\parindent}{0em}
\setlength{\parskip}{0.6em}
\renewcommand{\baselinestretch}{1.0}
\begin{document}
\begin{center}
{\LARGE \textit{Baker v. General Motors Corp.}, 522 U.S. 222 (1998)\par}
\vspace{0.5em}
{\large Case Brief --- Two-Page Edition\par}
\end{center}

\section{Memory Jogger}
Sister-state judgments get exacting full faith and credit with no public-policy exception, but a Michigan gag injunction could not stop a former GM engineer from testifying under subpoena in a Missouri suit by strangers to the decree (\href{https://supreme.justia.com/cases/federal/us/522/222/}{522 U.S. 222}).

\section{Detailed Case Facts}
Former GM engineer Ronald Elwell settled a wrongful-discharge suit in Michigan; the court entered a stipulated injunction barring him from testifying against GM without consent. Beverly Garner died in a 1990 Missouri crash and fire in a Chevrolet Blazer, and her sons sued GM and subpoenaed Elwell.

\section{Procedural History}
The federal district court in Missouri let Elwell testify, and the jury awarded \$11.3 million. The Eighth Circuit reversed on full faith and credit grounds, and the Supreme Court granted certiorari.

\section{Judicial Votes}
Reversed, 9--0 in the judgment: Ginsburg wrote for five; Scalia and Kennedy (with O'Connor and Thomas) concurred in the judgment.

\section{Holding}
Full faith and credit for final judgments is exacting and has no public-policy exception, but enforcement measures remain governed by forum law, and a sister-state decree cannot control witness testimony or subpoenas in litigation by strangers to it.

\section{Analysis of Opinions}
Ginsburg distinguished laws, which a forum may decline to apply, from judgments, which bind nationwide, as in \textit{Fauntleroy v.\ Lum}. Equitable decrees are not exempt, but the forum need not adopt the rendering state's enforcement mechanisms, and Michigan had no authority over the Bakers. The decree's Georgia carve-out and GM's side agreement showed GM was using Missouri as a ``surrogate tribunal.'' Scalia would hold that judgments cannot be executed out of state. Kennedy would ask only what preclusive effect Michigan itself would give, and Michigan does not bind nonparties.

\section{Examples: Future Applications}
\textbf{Same-Side:} nonparty plaintiffs in State~B may subpoena a scientist gagged by a State~A decree. \textbf{Opposite-Side:} a State~A money judgment for breaching a confidentiality settlement must be enforced in State~B despite contrary policy. \textbf{Fence-Sitter:} a State~A judgment against a shield-law physician for telehealth abortion prescriptions.

\section{Critique}
Kennedy warned that exempting injunctions that interfere with sister-state litigation smuggles in a public-policy exception under a jurisdictional label. Ginsburg replied that mutuality doctrine alone cannot explain why Missouri may compel Elwell, a Michigan party, to defy a valid order.

\section{Key Quotations}
\begin{quote}
``Regarding judgments, however, the full faith and credit obligation is exacting. A final judgment in one State, if rendered by a court with adjudicatory authority over the subject matter and persons governed by the judgment, qualifies for recognition throughout the land. \ldots But our decisions support no roving `public policy exception' to the full faith and credit due judgments.''
\end{quote}
\noindent --- Justice Ginsburg, Opinion of the Court, 522 U.S. at 233.

\end{document}
